\begin{pastexam}[label={px:ncu-114-dsa-02}]{中央資工所 114 學年度 第 2 題}
（複選；每題 5 分，全部答對才給分，答錯倒扣 1 分）

Which of the following statements are correct?
\begin{enumerate}[label=(\alph*)]
\item The postfix of (a+b)/(c-d)+e is ab+cd-/+e.
\item The postfix of infix expression b+a/c*d is bac+d/*.
\item The infix of the postfix expression ab+c*d/ef-/ is (((a+b)*c)/d)/(e-f).
\item The prefix of the postfix expression abc+d/* is *a/+bcd.
\end{enumerate}
\end{pastexam}

\begin{pastsoln}
答案：(c)(d)。

中序轉後序照資結書第~6.3.3~節：括號和優先順序決定運算的先後，同優先順序由左往右結合；後序式
裡每個運算子緊接在它的兩個運算元後面。 後序轉回中序或前序，用一個堆疊由左往右掃（資結書
第~6.3.5~節）：遇到運算元就推入；遇到運算子就彈出兩個（先彈出的是右運算元），組合後推回去。

\textbf{(a) 錯}。 $(a+b)/(c-d)+e$ 的運算順序：$a+b$、$c-d$、兩者相除、最後加 $e$。 最後一個運算是
$+$，所以後序要以 $+$ 結尾：
\[
(a+b) \to \texttt{ab+},\quad (c-d) \to \texttt{cd-},\quad \div \to \texttt{ab+cd-/},\quad
+e \to \texttt{ab+cd-/e+}.
\]
題目的 \texttt{ab+cd-/+e} 把 $+$ 放在 $e$ 前面，$+$ 的運算元變成 \texttt{ab+cd-/} 底下的東西，
而且最後剩下 $e$ 沒有運算子可用，不是合法的後序式。

\textbf{(b) 錯}。 $b+a/c*d$：$/$ 和 $*$ 優先於 $+$，而 $/$、$*$ 同級，由左往右，所以是
$b + ((a/c)*d)$：
\[
a/c \to \texttt{ac/},\quad *d \to \texttt{ac/d*},\quad b+ \to \texttt{bac/d*+}.
\]
題目的 \texttt{bac+d/*} 對應的是 $b*((a+c)/d)$，運算子全錯。

\textbf{(c) 對}。 掃 \texttt{ab+c*d/ef-/}：
\begin{center}
\small
\begin{tabular}{cl}
\toprule
讀到 & 堆疊（底 $\to$ 頂） \\
\midrule
\texttt{a}、\texttt{b} & $a$, $b$ \\
\texttt{+} & $(a+b)$ \\
\texttt{c}、\texttt{*} & $((a+b)*c)$ \\
\texttt{d}、\texttt{/} & $(((a+b)*c)/d)$ \\
\texttt{e}、\texttt{f}、\texttt{-} & $(((a+b)*c)/d)$, $(e-f)$ \\
\texttt{/} & $(((a+b)*c)/d)/(e-f)$ \\
\bottomrule
\end{tabular}
\end{center}
和題目寫的一樣。

\textbf{(d) 對}。 掃 \texttt{abc+d/*}，這次組合成前序（運算子放前面）：\texttt{b}、\texttt{c}、
\texttt{+} 得 \texttt{+bc}；\texttt{d}、\texttt{/} 得 \texttt{/+bcd}；最後 \texttt{*} 的左運算元是
\texttt{a}、右運算元是 \texttt{/+bcd}，得 \texttt{*a/+bcd}。 對應的中序是 $a*((b+c)/d)$。 和題目一樣。
\end{pastsoln}
